\documentclass[acmsmall,screen,nonacm]{acmart}

\usepackage{amsmath,amssymb}
\usepackage{tikz-cd}
\usepackage{listings}
\usepackage{cleveref}

\lstdefinestyle{ari}{
  basicstyle=\ttfamily\small,
  commentstyle=\itshape\color{gray},
  keywordstyle=\bfseries,
  columns=fullflexible,
  keepspaces=true,
  literate={->}{$\to$}2,
  frame=single,
  xleftmargin=1em,
  xrightmargin=1em,
}

\lstdefinestyle{ali}{
  basicstyle=\ttfamily\small,
  commentstyle=\itshape\color{gray},
  columns=fullflexible,
  keepspaces=true,
  literate={->}{$\to$}2 {<<=}{$\ll\!=$}2 {\#0}{$\otimes$}2,
  frame=single,
  xleftmargin=1em,
  xrightmargin=1em,
  moredelim=[s]{(*}{*)},
  morecomment=[s]{(*}{*)},
}

\newcommand{\ob}{\mathsf{ob}}
\newcommand{\unit}{\mathsf{unit}}
\newcommand{\id}{\mathrm{id}}
\newcommand{\copy}{\mathsf{copy}}
\newcommand{\swap}{\mathsf{swap}}
\newcommand{\erase}{\mathsf{erase}}
\newcommand{\alifib}{\textsf{alifib}}

\begin{document}

\title{Translating Term Rewrite Systems to Directed Complexes}
\author{} \date{}
\maketitle

\begin{abstract}
We describe a tool that translates term rewrite systems (TRS) in ARI
format into \alifib{} directed complexes, using the standard encoding
from polygraphic rewriting theory.  Function symbols become
$2$\nobreakdash-cells, rewrite rules become $3$\nobreakdash-cells, and
structural cells (copying, permutation, erasure) mediate between the
linear world of string diagrams and the non-linear world of terms.
\end{abstract}

\tableofcontents

% ===================================================================
\section{Usage}
\label{sec:usage}
% ===================================================================

The converter lives in \texttt{trs/} and is written in TypeScript.
Build and run as follows.

\medskip\noindent\textbf{Build.}
\begin{verbatim}
  cd trs/
  npx tsc          # compiles to dist/
\end{verbatim}

\noindent\textbf{Run.}
\begin{verbatim}
  node dist/main.js <input.ari> [output.ali]
\end{verbatim}
If no output path is given, the tool writes to the input path with
extension changed to \texttt{.ali}.  Use \texttt{-} as the output to
print to stdout.

\medskip\noindent\textbf{Verify.}
Feed the output to \alifib:
\begin{verbatim}
  alifib output.ali
\end{verbatim}

\noindent The tool reads the ARI format used by TPDB: S-expression
declarations \texttt{(format TRS)}, \texttt{(fun $f$ $n$)}, and
\texttt{(rule $\ell$ $r$)}.  Quoted symbols \texttt{|0|} are unquoted
and sanitized into valid \alifib{} identifiers.

% ===================================================================
\section{Example}
\label{sec:example}
% ===================================================================

Consider the following TRS, expressing that a unary operation
$\mathsf{minus}$ is an involutive anti-homomorphism with respect to a
unary $h$ and a binary $f$.

\begin{lstlisting}[style=ari,title={Input: \texttt{involution.ari}}]
(format TRS)
(fun minus 1)
(fun h 1)
(fun f 2)
(rule (minus (minus x)) x)
(rule (minus (h x)) (h (minus x)))
(rule (minus (f x y)) (f (minus y) (minus x)))
\end{lstlisting}

\noindent Running the converter produces the following \alifib{} file.

\begin{lstlisting}[style=ali,title={Output: \texttt{involution.ali}}]
@Type
Eq <<= {
  pt,
  dom : pt -> pt, cod : pt -> pt,
  lhs : dom -> cod, rhs : dom -> cod,
  dir : lhs -> rhs, inv : rhs -> lhs
},

involution <<= {
  pt,
  ob : pt -> pt,

  (* Structural cells *)
  copy : ob -> ob ob,
  swap : ob ob -> ob ob,

  (* Function symbols *)
  minus : ob -> ob,
  h     : ob -> ob,
  f     : ob ob -> ob,

  (* Identity 2-cells *)
  id_1 : ob -> ob,
  id_2 : ob ob -> ob ob,

  (* Naturality of copy *)
  attach Copy_minus :: Eq along
    [ lhs => minus copy,
      rhs => copy (minus #0 minus) ],
  attach Copy_h :: Eq along
    [ lhs => h copy,
      rhs => copy (h #0 h) ],
  attach Copy_f :: Eq along
    [ lhs => f copy,
      rhs => (copy #0 copy) (ob swap ob) (f #0 f) ],

  (* Structural equations *)
  id_1_idem     : id_1 id_1 -> id_1,
  id_1_idem_inv : id_1 -> id_1 id_1,
  swap_inv      : swap swap -> id_2,
  swap_inv_inv  : id_2 -> swap swap,
  copy_comm     : copy swap -> copy,
  copy_comm_inv : copy -> copy swap,

  (* Rewrite rules *)
  rule_1 : minus minus -> id_1,
  rule_2 : h minus -> minus h,
  rule_3 : f minus -> swap (minus #0 minus) f
}
\end{lstlisting}

\noindent The three rewrite rules translate to $3$-cells whose
source and target are the two sides of each rule, encoded as
$2$-morphisms in the polygraph.

% ===================================================================
\section{The polygraphic encoding}
\label{sec:encoding}
% ===================================================================

The translation from a TRS to a directed complex follows the standard
theory of polygraphic rewriting~\cite{Burroni1993,Guiraud2019}.  We
briefly recall the key ideas and then describe the details of the
encoding.

% -------------------------------------------------------------------
\subsection{Background: polygraphs}
\label{sec:background}
% -------------------------------------------------------------------

A \emph{polygraph} (also called a \emph{computad}) is a
higher-dimensional presentation of an algebraic structure.  Just as a
group presentation has generators and relations, a polygraph has
generators at each dimension:
\begin{itemize}
\item \textbf{$0$-cells} are points (objects).
\item \textbf{$1$-cells} are arrows between points.
\item \textbf{$2$-cells} are rewrites between (composites of) arrows.
\item \textbf{$3$-cells} are rewrites between $2$-cell composites.
\end{itemize}
In the context of term rewriting, a TRS with signature $\Sigma$ and
rules $R$ gives rise to a $(3,2)$-polygraph: function symbols are
$2$-generators and rewrite rules are $3$-generators.

The key insight is that terms are \emph{tree-shaped} (they can share
and discard variables), while cells in a polygraph are
\emph{string-diagram-shaped} (every input wire goes to exactly one
output).  To bridge this gap, we introduce \emph{structural cells} that
mediate between the linear world of string diagrams and the non-linear
world of terms.

% -------------------------------------------------------------------
\subsection{Dimension 0: the point}
\label{sec:dim0}
% -------------------------------------------------------------------

Every complex begins with a single $0$-cell:
\[
  \mathsf{pt}
\]
This is the unique object of the underlying monoid (or, if one prefers,
the unique $0$-cell of the one-object free higher category).

% -------------------------------------------------------------------
\subsection{Dimension 1: the sort}
\label{sec:dim1}
% -------------------------------------------------------------------

For an unsorted (single-sorted) TRS, we introduce one $1$-cell:
\[
  \ob : \mathsf{pt} \to \mathsf{pt}
\]
This plays the role of the unique sort.  A term of the TRS with $n$
free variables will be encoded as a $2$-morphism
$\ob^n \to \ob$, where $\ob^n$ denotes the $n$-fold horizontal
composite $\ob \otimes \cdots \otimes \ob$.

When the TRS has \emph{constants} (function symbols of arity~$0$) or
\emph{erasing rules} (rules where a variable appears on the left but
not on the right), we also need a second $1$-cell:
\[
  \unit : \mathsf{pt} \to \mathsf{pt}
\]
This is the monoidal unit stand-in.  A constant $c$ is then a $2$-cell
$\unit \to \ob$ rather than $\to \ob$ (which is not syntactically
well-formed as a boundary in \alifib).

% -------------------------------------------------------------------
\subsection{Dimension 2: function symbols and structural cells}
\label{sec:dim2}
% -------------------------------------------------------------------

\subsubsection{Function symbols.}
Each function symbol $f$ of arity~$n$ becomes a $2$-cell:
\[
  f : \underbrace{\ob\;\ob\;\cdots\;\ob}_{n} \to \ob
\]
For a constant $c$ of arity~$0$:
\[
  c : \unit \to \ob
\]
A term $f(t_1, \ldots, t_n)$ is then encoded by first encoding each
$t_i$ as a $2$-morphism (recursively), composing them horizontally
with $\otimes$, and then composing vertically with $f$.  For example,
$h(\mathsf{minus}(x))$ is encoded as
$\mathsf{minus} \cdot h$
(apply $\mathsf{minus}$, then $h$), and $f(t_1, t_2)$ is encoded as
$(t_1 \otimes t_2) \cdot f$ where $\otimes$ is horizontal composition
(\texttt{\#0} in \alifib).

\subsubsection{Structural cells.}
String diagrams are \emph{linear}: each input wire is used exactly
once.  Terms are not.  To handle variable sharing (contraction) and
discarding (weakening), we introduce:
\begin{align*}
  \copy  &: \ob \to \ob\;\ob  & &\text{(duplicate a wire)} \\
  \swap  &: \ob\;\ob \to \ob\;\ob & &\text{(exchange two wires)} \\
  \erase &: \ob \to \unit     & &\text{(discard a wire)}
\end{align*}
The $\erase$ cell is only generated when the TRS has constants or
erasing rules.  Together with the monoidal unit eliminators
$\mathsf{unit\_l} : \unit\;\ob \to \ob$ and
$\mathsf{unit\_r} : \ob\;\unit \to \ob$, these cells allow us to
convert any term into a string diagram over the signature.

\subsubsection{Identity cells.}
In \alifib, the horizontal composite $\id_\ob \otimes \id_\ob$
(written \texttt{id \#0 id}) produces a non-globular diagram that
cannot serve as the boundary of a higher cell.  We therefore introduce
explicit identity $2$-cells:
\[
  \id_n : \ob^n \to \ob^n
  \qquad\text{for each $n$ up to the maximum arity}
\]
In practice, $\id_1$ and $\id_2$ are always generated, and higher
$\id_n$ are added as needed.

% -------------------------------------------------------------------
\subsection{Encoding a term as a 2-morphism}
\label{sec:term-encoding}
% -------------------------------------------------------------------

Given a rule $\ell \to r$ with variables $x_1, \ldots, x_k$ (those
appearing in $\ell$, in left-to-right first-occurrence order), each
side is encoded as a $2$-morphism $\ob^k \to \ob$.  The encoding of a
term~$t$ starting from input wires $x_1, \ldots, x_k$ proceeds in two
phases.

\medskip\noindent\textbf{Phase 1: Gather.}
The \emph{gather phase} routes the $k$ input wires to match the leaf
variable pattern of $t$.  This involves three sub-steps:
\begin{enumerate}
\item \textbf{Copy:} For each variable $x_i$ that appears $m$ times in
  $t$, apply a copy tree $\ob \to \ob^m$.  For a variable that appears
  $0$ times, apply $\erase$.  These operations are composed
  horizontally across all input wires.

\item \textbf{Eliminate unit wires:} If any variables were erased, the
  resulting $\unit$ wires are removed using $\mathsf{unit\_l}$ and
  $\mathsf{unit\_r}$.

\item \textbf{Permute:} The copied wires may not be in the order
  required by $t$.  A \emph{swap network} of adjacent transpositions
  (using the $\swap$ cell) permutes them into the correct leaf order.
  The algorithm computes a bubble-sort sequence from the current wire
  arrangement to the target arrangement.
\end{enumerate}

\medskip\noindent\textbf{Phase 2: Apply.}
The \emph{apply phase} recursively encodes the term structure.  Each
subterm $f(t_1, \ldots, t_n)$ is encoded as:
\[
  (t_1 \otimes \cdots \otimes t_n) \cdot f
\]
where each $t_i$ is encoded recursively.  A variable leaf becomes the
identity~$\id_1$, and a constant leaf becomes the constant cell
directly.

\medskip\noindent\textbf{Example.}
Consider the rule $f(x, y) \to f(y, x)$ (commutativity).  Both sides
have variables $[x, y]$.
\begin{itemize}
\item \emph{LHS:} The leaf order of $f(x,y)$ is $[x,y]$, which matches
  the input order.  The gather phase is trivial.  The apply phase gives
  $f$.  So the LHS encoding is $f$.
\item \emph{RHS:} The leaf order of $f(y,x)$ is $[y,x]$.  The gather
  phase produces $\swap$ (one adjacent transposition).  The apply phase
  gives $f$.  So the RHS encoding is $\swap \cdot f$.
\end{itemize}
The rewrite rule becomes the $3$-cell $\mathsf{rule} : f \to \swap
\cdot f$.

\medskip\noindent\textbf{Example (non-linear).}
Consider $f(x, x) \to x$ (idempotence of $f$).  Variables: $[x]$.
\begin{itemize}
\item \emph{LHS:} The leaf pattern is $[x, x]$.  The gather phase
  produces $\copy$.  The apply phase gives $f$.  So the LHS encoding is
  $\copy \cdot f$.
\item \emph{RHS:} The leaf pattern is $[x]$.  The gather phase is
  trivial.  The apply phase gives $\id_1$.  So the RHS encoding is
  $\id_1$.
\end{itemize}
The rewrite rule becomes $\mathsf{rule} : \copy \cdot f \to \id_1$.

% -------------------------------------------------------------------
\subsection{Handling constants}
\label{sec:constants}
% -------------------------------------------------------------------

When a term contains a constant $c$ (arity $0$), the encoding needs to
produce a $\unit$ wire and then apply $c : \unit \to \ob$.  Since the
input context consists entirely of $\ob$ wires, the strategy is:
\begin{enumerate}
\item Make extra copies of a ``donor'' variable (the first input
  variable) to create enough $\ob$ wires for both variable slots and
  constant slots.
\item Permute all wires (while they are all $\ob$, so $\swap$ applies)
  to place the donor copies at the positions where constants will go.
\item Apply $\erase : \ob \to \unit$ at constant positions to convert
  $\ob$ wires to $\unit$ wires.
\item Apply the raw term encoding, which expects $\unit$ wires at
  constant leaf positions.
\end{enumerate}

% -------------------------------------------------------------------
\subsection{Dimension 3: rewrite rules and structural equations}
\label{sec:dim3}
% -------------------------------------------------------------------

\subsubsection{Rewrite rules.}
Each TRS rule $\ell \to r$ becomes a $3$-cell:
\[
  \mathsf{rule}_i : \llbracket \ell \rrbracket \to \llbracket r \rrbracket
\]
where $\llbracket t \rrbracket$ denotes the $2$-morphism encoding of
$t$.  Both sides have the same source ($\ob^k$) and target ($\ob$),
so the $3$-cell is well-formed.

\subsubsection{Naturality of copy.}
For the structural cells to interact correctly with function symbols, we
need \emph{naturality} equations: the $\copy$ cell commutes with
function application.  For a unary $f$:
\[
  f \cdot \copy = \copy \cdot (f \otimes f)
\]
For a binary $f$:
\[
  f \cdot \copy = (\copy \otimes \copy) \cdot
  (\ob\;\swap\;\ob) \cdot (f \otimes f)
\]
The right-hand side of the binary case uses a permutation to route the
two copies of the first argument together and the two copies of the
second argument together.

These are declared as \emph{Eq attachments} in \alifib{}: the type
$\mathsf{Eq}$ is a small complex with cells \texttt{pt}, \texttt{dom},
\texttt{cod}, \texttt{lhs}, \texttt{rhs}, \texttt{dir}, \texttt{inv},
and attaching along it produces a pair of invertible $3$-cells
(one in each direction) witnessing the equation.

\subsubsection{Structural equations.}
We also declare a handful of coherence cells:
\begin{align*}
  &\id_1 \cdot \id_1 \leftrightarrow \id_1
    &&\text{(identity is idempotent)} \\
  &\swap \cdot \swap \leftrightarrow \id_2
    &&\text{(swap is an involution)} \\
  &\copy \cdot \swap \leftrightarrow \copy
    &&\text{(copy is commutative)}
\end{align*}
Each is given as a pair of inverse $3$-cells.

% -------------------------------------------------------------------
\subsection{Summary of the encoding}
\label{sec:summary}
% -------------------------------------------------------------------

Given a TRS $(\Sigma, R)$, the output directed complex has:

\begin{center}
\begin{tabular}{cll}
\textbf{Dim} & \textbf{Generators} & \textbf{Source} \\
\hline
0 & $\mathsf{pt}$ & --- \\
1 & $\ob$, (optionally $\unit$) & sort(s) \\
2 & $f$ for each $f \in \Sigma$ & function symbols \\
  & $\copy$, $\swap$, ($\erase$, $\mathsf{unit\_l}$,
    $\mathsf{unit\_r}$) & structural cells \\
  & $\id_1, \id_2, \ldots$ & identity cells \\
3 & $\mathsf{rule}_i$ for each rule in $R$ & rewrite rules \\
  & $\mathsf{Copy\_}f.\mathsf{dir}/\mathsf{inv}$ for each $f$ & naturality \\
  & involutions for $\swap$, $\id_1$, $\copy$ & coherence
\end{tabular}
\end{center}

The resulting complex is a finite presentation of the free
$(\infty, 2)$-category generated by the TRS: the $2$-cells present the
term algebra, and the $3$-cells present the rewriting relation.  The
higher cells (naturality, coherence) ensure that the structural
operations interact with function symbols in the expected way.

Paths of $3$-cells in this complex correspond exactly to rewrite
sequences in the original TRS, modulo the structural rearrangements
needed to match variables in consecutive rule applications.

% ===================================================================

\begin{thebibliography}{9}

\bibitem{Burroni1993}
A.~Burroni.
\newblock Higher-dimensional word problems with applications to equational
  logic.
\newblock \emph{Theoretical Computer Science}, 115(1):43--62, 1993.

\bibitem{Guiraud2019}
Y.~Guiraud and P.~Malbos.
\newblock Polygraphs of finite derivation type.
\newblock \emph{Mathematical Structures in Computer Science},
  29(8):1130--1192, 2019.

\end{thebibliography}

\end{document}
